% ============================================================
%  Übungsaufgaben Template (my_dhbw Stil, uebung_*.tex)
%  Header "Übungsblatt", grün eingefärbte <details>-Lösungen.
%  Renderer: pdflatex (zweimal)
% ============================================================
\documentclass[11pt, a4paper]{article}

\usepackage[ngerman]{babel}
\usepackage{fontspec}
\setmainfont[
  Path=/usr/share/fonts/TTF/,
  Extension=.ttf,
  BoldFont=DejaVuSans-Bold,
  ItalicFont=DejaVuSans-Oblique,
  BoldItalicFont=DejaVuSans-BoldOblique
]{DejaVuSans}
\setsansfont[
  Path=/usr/share/fonts/TTF/,
  Extension=.ttf,
  BoldFont=DejaVuSans-Bold,
  ItalicFont=DejaVuSans-Oblique,
  BoldItalicFont=DejaVuSans-BoldOblique
]{DejaVuSans}
\setmonofont[Path=/usr/share/fonts/TTF/,Extension=.ttf]{DejaVuSansMono}
\usepackage[top=2cm, bottom=2cm, left=2.4cm, right=2.4cm, footskip=1cm]{geometry}
\usepackage{amsmath, amssymb}
\usepackage{xcolor}
\usepackage{enumitem}
\usepackage{fancyhdr}
\usepackage{lastpage}
\usepackage{tabularx}
\usepackage{booktabs}
\usepackage{microtype}
\usepackage{parskip}

\usepackage{array}
\usepackage{longtable}
\usepackage{ltablex}
\keepXColumns
\usepackage{colortbl}
\usepackage[most,breakable]{tcolorbox}
\usepackage{listings}
\usepackage[normalem]{ulem}
\usepackage{pdflscape}
\usepackage{titlesec}
\usepackage[colorlinks=true, linkcolor=accent, urlcolor=accent]{hyperref}

\renewcommand{\familydefault}{\sfdefault}

% ── Theme ─────────────────────────────────────────────────────
\definecolor{accent}{RGB}{0, 60, 113}
\definecolor{primaryblue}{RGB}{0, 60, 113}
\definecolor{loesung}{RGB}{0, 100, 0}
\definecolor{codebg}{RGB}{245, 245, 245}
\definecolor{figurebg}{RGB}{232, 241, 255}
\definecolor{tablehintbg}{RGB}{240, 255, 240}
\definecolor{solutionbg}{RGB}{240, 252, 240}
\definecolor{rulegray}{RGB}{180, 180, 180}
\definecolor{tableheadbg}{RGB}{0, 60, 113}

% ── Header/Footer ─────────────────────────────────────────────
\pagestyle{fancy}
\fancyhf{}
\renewcommand{\headrulewidth}{0.4pt}
\fancyhead[L]{\footnotesize\color{accent}\nichttitle}
\fancyhead[R]{\footnotesize\color{accent}\nichtsubtitle}
\fancyfoot[R]{\footnotesize\color{accent}Blatt \thepage\,/\,\pageref{LastPage}}

% ── Typografie ────────────────────────────────────────────────
\titleformat{\section}
    {\color{accent}\large\bfseries}{}{0pt}{}
    [\vspace{-4pt}\color{accent}\rule{\linewidth}{1.5pt}]
\titleformat{\subsection}
    {\normalsize\bfseries\color{accent!85!black}}{}{0pt}{}{}
\titleformat{\subsubsection}
    {\small\bfseries\color{accent!70!black}}{}{0pt}{}{}
\titlespacing*{\section}{0pt}{10pt}{3pt}
\titlespacing*{\subsection}{0pt}{6pt}{2pt}
\titlespacing*{\subsubsection}{0pt}{4pt}{1pt}

\setlength{\parindent}{0pt}
\setlength{\parskip}{6pt}
\renewcommand{\arraystretch}{1.2}
\setlist[itemize]{noitemsep, topsep=3pt, parsep=0pt, leftmargin=1.4em}
\setlist[enumerate]{noitemsep, topsep=3pt, parsep=0pt, leftmargin=1.6em}

% ── Boxen: solutionbox = Lösungsbox in grün ──────────────────
\newtcolorbox{figurebox}[1]{
  colback=figurebg, colframe=accent!50,
  title={Abbildung: #1}, fonttitle=\small\bfseries,
  breakable, before skip=6pt, after skip=6pt
}
\newtcolorbox{tablehintbox}[1]{
  colback=tablehintbg, colframe=green!50!black!50,
  title={Tabelle: #1}, fonttitle=\small\bfseries,
  breakable, before skip=6pt, after skip=6pt
}
\newtcolorbox{solutionbox}[1]{
  enhanced,
  colback=solutionbg, colframe=loesung,
  title={#1}, fonttitle=\small\bfseries,
  coltitle=white, colbacktitle=loesung,
  breakable, before skip=8pt, after skip=8pt
}

\lstset{
  basicstyle=\ttfamily\small,
  backgroundcolor=\color{codebg},
  breaklines=true,
  frame=single, framerule=0pt,
  xleftmargin=4pt, xrightmargin=4pt,
  aboveskip=6pt, belowskip=6pt,
  keepspaces=true
}

\providecommand{\nichttitle}{Übungsblatt}
\providecommand{\nichtsubtitle}{}

\renewcommand{\nichttitle}{Betriebssysteme}
\renewcommand{\nichtsubtitle}{Übungsblatt 9}


\begin{document}

\begin{center}
{\Large\bfseries\color{accent}\nichttitle}
\ifx\nichtsubtitle\empty\else
  \\[2pt]{\normalsize\color{accent!80!black}\nichtsubtitle}
\fi
\\[2pt]
\color{accent}\rule{0.4\linewidth}{0.8pt}\color{black}
\end{center}
\vspace{4pt}

% ── BODY ──────────────────────────────────────────────────────

\section{Übungsblatt 9 — Synchronisation}


\noindent\textcolor{rulegray}{\rule{\linewidth}{0.4pt}}


\paragraph{Aufgabe 1 — Race Condition im Kontostand \textit{(12 Punkte)}}\mbox{}\\[2pt]

Zwei Threads \texttt{T1} und \texttt{T2} bearbeiten ein gemeinsames Konto mit Startwert \texttt{saldo = 100}. Beide führen folgenden Pseudocode aus:


\begin{lstlisting}
1: tmp = saldo
2: tmp = tmp + 50
3: saldo = tmp
\end{lstlisting}


a) Definiere \textbf{Race Condition} und \textbf{Lost Update}. Nenne die zwei Voraussetzungen, unter denen eine Race Condition überhaupt entstehen kann. \textit{(3 P)}


b) Konstruiere eine Verschränkung der Befehle 1–3 von \texttt{T1} und \texttt{T2}, bei der am Ende \texttt{saldo = 150} ist (statt der erwarteten 200). Stelle den Ablauf in einer Tabelle mit den Spalten \textit{Schritt | T1 | T2 | tmp\_T1 | tmp\_T2 | saldo} dar. \textit{(5 P)}


c) Begründe, warum das Sperren der Interrupts auf einem Einkernsystem das Problem löst, auf einem Mehrkernsystem aber nicht ausreicht. \textit{(4 P)}


\textbf{Lösung:}


a) Eine \textbf{Race Condition} liegt vor, wenn das Ergebnis einer Berechnung von der zeitlichen Reihenfolge der Zugriffe nebenläufiger Aktivitäten abhängt. Voraussetzungen: (1) \textbf{Nebenläufigkeit} (mind. zwei quasi- oder echt-parallel laufende Aktivitäten) und (2) \textbf{gemeinsame Betriebsmittel} (hier \texttt{saldo}). Ein \textbf{Lost Update} ist der Spezialfall, dass Prozess A zuerst liest, B liest denselben alten Wert, beide schreiben zurück und das Update von A überschrieben (verloren) wird.


b) Verschränkung mit Lost Update:



{\small
\begin{tabularx}{\linewidth}{XXXXXX}
\hline
\rowcolor{tableheadbg}
{\color{white}\textbf{Schritt}} & {\color{white}\textbf{T1}} & {\color{white}\textbf{T2}} & {\color{white}\textbf{tmp\_T1}} & {\color{white}\textbf{tmp\_T2}} & {\color{white}\textbf{saldo}} \\[2pt]
\hline
\endhead
\hline
\endfoot
0 & – & – & – & – & 100 \\
1 & tmp = saldo & – & 100 & – & 100 \\
2 & tmp = tmp+50 & – & 150 & – & 100 \\
3 & – & tmp = saldo & 150 & 100 & 100 \\
4 & saldo = tmp & – & 150 & 100 & 150 \\
5 & – & tmp = tmp+50 & 150 & 150 & 150 \\
6 & – & saldo = tmp & 150 & 150 & \textbf{150} \\
\end{tabularx}
}


T2 hat den vor seiner Schreiboperation aktuellen Wert (150) überschrieben mit seinem auf veraltetem Lesen basierenden Ergebnis (150) — das Update von T1 ist verloren.


c) Auf einem \textbf{Einkernsystem} wird Nebenläufigkeit ausschließlich durch Interrupts (Timer-Interrupt für Kontextwechsel) erzeugt. Maskiert man Interrupts während der drei Befehle, kann kein anderer Thread dazwischen zur CPU kommen — der Block wird atomar. Auf einem \textbf{Mehrkernsystem} läuft \texttt{T2} jedoch \textit{echt-parallel} auf einem zweiten Kern; das Maskieren der Interrupts auf dem eigenen Kern hindert den anderen Kern nicht am Zugriff auf \texttt{saldo}. Das Race Condition bleibt bestehen — daher braucht man hier hardwaregestützte atomare Befehle (TSL/TAS), die den Speicherbus während Test+Set blockieren.



\noindent\textcolor{rulegray}{\rule{\linewidth}{0.4pt}}


\paragraph{Aufgabe 2 — TSL-Spinlock korrekt einsetzen \textit{(15 Punkte)}}\mbox{}\\[2pt]

Gegeben sei der Maschinenbefehl \texttt{TSL reg, addr}, der den Wert an \texttt{addr} in \texttt{reg} lädt und gleichzeitig (im selben Buszyklus) \texttt{addr := 1} setzt.


Ein Studierender schlägt folgenden Code für \texttt{enter\_region} und \texttt{leave\_region} vor:


\begin{lstlisting}
enter_region:
    MOV reg, #1
    XCHG reg, lock     ; tauscht reg <-> lock atomar
    CMP reg, #0
    JNE enter_region
    RET

leave_region:
    MOV lock, #0
    RET
\end{lstlisting}


a) Erkläre, warum der Befehl TSL für die Mutual-Exclusion-Bedingung \textit{zwingend} in einer einzigen, ununterbrechbaren Operation ausgeführt werden muss. Was würde passieren, wenn er aus zwei separaten Maschinenbefehlen (Lesen, Schreiben) bestünde? \textit{(4 P)}


b) Prüfe, ob der oben angegebene Code korrekt ist. Argumentiere für jede der vier Mutex-Bedingungen (Mutual Exclusion, Unabhängigkeit von Geschwindigkeit/CPU-Anzahl, kein Blockieren außerhalb, Fairness) ob sie erfüllt ist. \textit{(7 P)}


c) Nenne \textit{einen} konkreten Anwendungsfall, in dem ein Spinlock einem Semaphor vorzuziehen ist, sowie \textit{einen}, in dem das Gegenteil gilt. Begründe jeweils anhand von Wartezeit und CPU-Auslastung. \textit{(4 P)}


\textbf{Lösung:}


a) Wäre TSL aus zwei Befehlen (Schritt 1: lesen, Schritt 2: schreiben) zusammengesetzt, könnte zwischen ihnen ein Kontextwechsel oder ein paralleler Zugriff eines anderen Kerns stattfinden. Beispiel: Prozess A liest \texttt{lock = 0}, bevor er schreiben kann unterbricht ihn Prozess B, der ebenfalls \texttt{lock = 0} liest, beide setzen anschließend \texttt{lock = 1} und beide betreten den kritischen Abschnitt — Mutual Exclusion ist verletzt. Daher muss Test+Set in \textit{einem Speicherzugriffszyklus} ablaufen; die CPU blockiert dafür den Bus, sodass kein anderer Kern in dieser Zeit auf \texttt{lock} zugreifen kann.


b) Bewertung der vier Bedingungen:


\begin{itemize}
  \item \textbf{Mutual Exclusion} ✓: \texttt{XCHG} ist atomar; nur der Prozess, der vom alten \texttt{lock} den Wert 0 lädt (also \texttt{reg = 0} nach dem Tausch), verlässt die Schleife. Alle anderen sehen \texttt{reg = 1} und drehen weiter.
  \item \textbf{Unabhängigkeit von Geschwindigkeit / Anzahl Prozessoren} ✓: TSL/XCHG nutzt eine Bus-Sperre und funktioniert auf Einkern wie auf Mehrkernsystemen, unabhängig von Taktfrequenzen.
  \item \textbf{Kein Blockieren außerhalb} ✓: Außerhalb des kritischen Abschnitts ist \texttt{lock = 0}, jeder Wartende kann eintreten.
  \item \textbf{Fairness} ✗: Es gibt keine Warteschlange. Wer den Bus zuerst „erwischt", bekommt den Lock — ein Prozess kann theoretisch beliebig oft erfolglos spinnen, während andere ihn überholen (\textbf{Starvation/Verhungern} möglich). Diese Bedingung ist \textit{nicht} erfüllt.
\end{itemize}

c) \textbf{Spinlock vorzuziehen}: kurze, vorhersehbare Wartezeit (z. B. ein paar Instruktionen im Kernel zur Aktualisierung einer Statistik) auf einem Mehrkernsystem — der Aufwand für Kontextwechsel + Aufwecken (Semaphor) wäre teurer als kurze Busy-Wait-Zyklen. \textbf{Semaphor vorzuziehen}: lange oder unbekannte Wartezeit (z. B. Warten auf I/O oder einen vollen Drucker-Puffer) — Spinning würde die CPU verschwenden; ein Semaphor suspendiert den Prozess und gibt die CPU frei.



\noindent\textcolor{rulegray}{\rule{\linewidth}{0.4pt}}


\paragraph{Aufgabe 3 — Erzeuger-Verbraucher mit Zählsemaphoren \textit{(20 Punkte)}}\mbox{}\\[2pt]

Ein Ringpuffer fasst \texttt{N = 4} Slots. Es laufen \texttt{2} Erzeuger und \texttt{3} Verbraucher nebenläufig. Jeder Erzeuger-Aufruf legt ein Element ab, jeder Verbraucher-Aufruf entnimmt eines.


a) Initialisiere drei Semaphore \texttt{mutex}, \texttt{empty}, \texttt{full} mit den korrekten Startwerten und erkläre die Rolle jedes einzelnen. \textit{(5 P)}


b) Schreibe den Pseudocode für \texttt{produce()} und \texttt{consume()} unter Verwendung der Semaphore. Achte auf die Reihenfolge der \texttt{P()}-Operationen. \textit{(7 P)}


c) Ein Studierender vertauscht in \texttt{produce()} die ersten beiden \texttt{P()}-Operationen: erst \texttt{P(mutex)}, dann \texttt{P(empty)}. Konstruiere ein Szenario, in dem dadurch eine \textbf{Verklemmung} entsteht. Stelle den Zustand der drei Zähler und die Wartelisten dar. \textit{(8 P)}


\textbf{Lösung:}


a) Initialisierung:


\begin{itemize}
  \item \texttt{mutex = 1} (binäres Semaphor) — schützt den kritischen Abschnitt um die Pufferdatenstruktur (Schreib-/Lesezeiger), damit nicht zwei Erzeuger gleichzeitig in dieselbe Zelle schreiben.
  \item \texttt{empty = 4} (Zählsemaphor) — zählt freie Slots. Erzeuger müssen \texttt{P(empty)} machen; sind keine frei (Zähler 0), werden sie suspendiert.
  \item \texttt{full = 0} (Zählsemaphor) — zählt belegte Slots. Verbraucher machen \texttt{P(full)}; ist nichts da, werden sie suspendiert.
\end{itemize}

b) Pseudocode:


\begin{lstlisting}
produce():
    P(empty)        // freien Slot reservieren
    P(mutex)        // KA betreten
    buffer[in] = item; in = (in+1) % N
    V(mutex)        // KA verlassen
    V(full)         // einen belegten Slot signalisieren

consume():
    P(full)
    P(mutex)
    item = buffer[out]; out = (out+1) % N
    V(mutex)
    V(full ? nein → empty)   // einen freien Slot freigeben
    V(empty)
\end{lstlisting}


Wichtig: Die Reihenfolge \texttt{P(empty)} \textit{vor} \texttt{P(mutex)} ist zwingend — andernfalls sperrt ein blockierender Erzeuger den Mutex und verklemmt alle Verbraucher.


c) Verklemmungsszenario bei vertauschter Reihenfolge in \texttt{produce()}:


\begin{lstlisting}
produce_falsch():
    P(mutex)        // <-- vertauscht
    P(empty)
    ...
\end{lstlisting}


Ablauf, Startzustand \texttt{mutex=1, empty=0, full=4} (Puffer komplett voll):



{\footnotesize
\begin{tabularx}{\linewidth}{XXXXXXX}
\hline
\rowcolor{tableheadbg}
{\color{white}\textbf{Schritt}} & {\color{white}\textbf{Aktion}} & {\color{white}\textbf{mutex}} & {\color{white}\textbf{empty}} & {\color{white}\textbf{full}} & {\color{white}\textbf{Wartend auf empty}} & {\color{white}\textbf{Wartend auf mutex}} \\[2pt]
\hline
\endhead
\hline
\endfoot
0 & Start (Puffer voll) & 1 & 0 & 4 & – & – \\
1 & Erzeuger E1 ruft \texttt{produce\_falsch()}: \texttt{P(mutex)} → mutex=0 & 0 & 0 & 4 & – & – \\
2 & E1: \texttt{P(empty)} → suspendiert (empty=0) & 0 & 0 & 4 & E1 & – \\
3 & Verbraucher V1 ruft \texttt{consume()}: \texttt{P(full)} → full=3 & 0 & 0 & 3 & E1 & – \\
4 & V1: \texttt{P(mutex)} → suspendiert (mutex=0) & 0 & 0 & 3 & E1 & V1 \\
5 & V2 analog & 0 & 0 & 2 & E1 & V1, V2 \\
6 & V3 analog & 0 & 0 & 1 & E1 & V1, V2, V3 \\
\end{tabularx}
}


\textbf{Verklemmung}: E1 hält \texttt{mutex} und wartet auf \texttt{empty}. Nur ein Verbraucher könnte \texttt{empty} erhöhen — alle Verbraucher hängen aber an \texttt{P(mutex)}, weil E1 ihn nie freigibt. Klassischer Deadlock: zyklische Wartebeziehung E1 → empty → V\_i → mutex → E1.


In der korrekten Variante (b) hätte E1 zuerst \texttt{P(empty)} ausgeführt, ohne \texttt{mutex} zu halten — die Verbraucher hätten den \texttt{mutex} problemlos bekommen, Slots geleert und \texttt{V(empty)} aufgerufen.



\noindent\textcolor{rulegray}{\rule{\linewidth}{0.4pt}}


\paragraph{Aufgabe 4 — Monitor: Mesa vs. Hoare \textit{(15 Punkte)}}\mbox{}\\[2pt]

Ein bounded buffer wird als Monitor implementiert. Beim Entnehmen wartet ein Konsument auf der Bedingung \texttt{notEmpty}:


\begin{lstlisting}
procedure get():
    if (count == 0) wait(notEmpty);
    item = buffer[out]; out = (out+1) % N; count--;
    signal(notFull);
    return item;
\end{lstlisting}


a) Erkläre die Aufgabe des „Türstehers" eines Monitors und worin der Unterschied zur Verwendung eines Semaphors besteht. \textit{(4 P)}


b) Ein Konsument K1 wartet in \texttt{wait(notEmpty)}. Produzent P1 legt ein Element ab und ruft \texttt{signal(notEmpty)}. Beschreibe den weiteren Ablauf jeweils unter \textbf{Signal-and-Continue (Mesa)} und unter \textbf{Signal-and-Wait (Hoare)}. \textit{(5 P)}


c) Begründe, warum unter Mesa-Semantik das \texttt{if} in \texttt{get()} durch ein \texttt{while} ersetzt werden \textbf{muss}, unter Hoare-Semantik aber \texttt{if} korrekt bleibt. Konstruiere ein Szenario mit einem zweiten Konsumenten K2, das die Notwendigkeit zeigt. \textit{(6 P)}


\textbf{Lösung:}


a) Der \textbf{Türsteher} stellt sicher, dass zu jedem Zeitpunkt höchstens ein Prozess innerhalb des Monitors aktiv ist — die Synchronisation wird vom Compiler automatisch um die Monitorprozeduren herum erzeugt, der Programmierer muss sie \textit{nicht} explizit codieren. Beim Semaphor hingegen ist der Programmierer selbst dafür verantwortlich, jedes \texttt{P()} mit dem passenden \texttt{V()} zu paaren — vergisst er ein \texttt{V()} oder verschachtelt fehlerhaft, droht Deadlock oder Lost Wakeup. Monitore verlagern Korrektheit von Laufzeit auf Übersetzungszeit.


b) Ablauf nach \texttt{signal(notEmpty)} durch P1:


\begin{itemize}
  \item \textbf{Mesa (Signal-and-Continue)}: P1 läuft im Monitor weiter, beendet seine Prozedur und verlässt den Monitor. K1 wird aus der Warteschlange auf \texttt{notEmpty} in die Eintritts-Warteschlange des Monitors verschoben, bekommt den Monitor irgendwann später und setzt seine Ausführung \textit{nach} \texttt{wait(notEmpty)} fort. Zwischen Signal und Wiederaufnahme können andere Monitor-Aufrufe stattfinden.
  \item \textbf{Hoare (Signal-and-Wait)}: P1 wird sofort suspendiert (und in eine spezielle „signaler"-Warteschlange gestellt), K1 erhält den Monitor unmittelbar nach \texttt{signal} und setzt nach \texttt{wait(notEmpty)} fort. Wenn K1 den Monitor verlässt, wird P1 reaktiviert.
\end{itemize}

c) Mesa-Szenario, das ein \texttt{while} erzwingt — Puffer leer, Konsumenten K1, K2 warten auf \texttt{notEmpty}:



{\small
\begin{tabularx}{\linewidth}{XXXX}
\hline
\rowcolor{tableheadbg}
{\color{white}\textbf{Schritt}} & {\color{white}\textbf{Akteur}} & {\color{white}\textbf{Aktion}} & {\color{white}\textbf{count}} \\[2pt]
\hline
\endhead
\hline
\endfoot
1 & P1 & betritt Monitor, legt Item ab, count=1 & 1 \\
2 & P1 & \texttt{signal(notEmpty)} — K1 wechselt in Eintritts-Queue (Mesa: P1 läuft weiter) & 1 \\
3 & P1 & verlässt Monitor & 1 \\
4 & K2 & tritt vor K1 in den Monitor ein, prüft count==0 → falsch, entnimmt Item, count=0 & 0 \\
5 & K2 & verlässt Monitor & 0 \\
6 & K1 & tritt ein, setzt nach \texttt{wait} fort — bei \texttt{if} läuft K1 ohne erneute Prüfung weiter und greift auf einen leeren Puffer zu! & 0 \\
\end{tabularx}
}


Mit \texttt{while (count == 0) wait(notEmpty);} würde K1 in Schritt 6 die Bedingung erneut testen, sehen, dass der Puffer wieder leer ist, und sich wieder in \texttt{wait} legen. Korrekt.


Unter \textbf{Hoare} kann zwischen \texttt{signal} und der Wiederaufnahme von K1 \textit{kein} anderer Prozess in den Monitor — K2 hätte keine Chance, sich vorzudrängen. Sobald K1 nach \texttt{wait} weiterläuft, ist der Zustand garantiert exakt der, den P1 hergestellt hat (\texttt{count == 1}). Daher genügt \texttt{if}. Die Schwäche von Mesa ist also der Preis für die einfachere Implementierung (kein Übergeben des Monitors, keine Signaler-Queue).



\noindent\textcolor{rulegray}{\rule{\linewidth}{0.4pt}}


\paragraph{Aufgabe 5 — Transfer: Lese-/Schreibsperre selbst gebaut \textit{(18 Punkte)}}\mbox{}\\[2pt]

Im Kapitel ist nur der einfache Mutex behandelt. Übertrage das Konzept auf eine \textbf{Reader-Writer-Sperre}: beliebig viele Leser dürfen gleichzeitig im kritischen Abschnitt sein, aber wenn ein Schreiber schreibt, ist niemand anderes drin.


a) Welcher Mutex-Bedingung (1–4 aus dem Kapitel) widerspricht es zunächst, mehrere Leser gleichzeitig in den KA zu lassen? Argumentiere, warum die Reader-Writer-Sperre dennoch \textit{nicht} die Korrektheit der Anwendung gefährdet. \textit{(4 P)}


b) Skizziere die Implementierung der Operationen \texttt{acquire\_read()}, \texttt{release\_read()}, \texttt{acquire\_write()}, \texttt{release\_write()} mithilfe eines Zählers \texttt{readers} (geschützt durch \texttt{mutex}) und eines binären Semaphors \texttt{wrt}. \textit{(8 P)}


c) Zeige, dass deine Lösung einen Schreiber \textbf{verhungern} lassen kann. Konstruiere eine Folge von Reader-/Writer-Aufrufen und zeige, welche Mutex-Bedingung verletzt wird. Wie würdest du das beheben (Skizze, kein Code)? \textit{(6 P)}


\textbf{Lösung:}


a) Es widerspricht der \textbf{Mutual-Exclusion}-Bedingung (Bedingung 1: „nie zwei gleichzeitig drin"). Die Anwendung bleibt aber korrekt, weil das Mutex-Modell ME zum Schutz \textit{gegen Race Conditions auf Schreibzugriffe} fordert. Race Conditions entstehen nur, wenn mindestens \textit{ein} Zugriff schreibt — gleichzeitige Lesezugriffe können einen konsistenten Zustand nicht zerstören. Die Bedingung wird also abgeschwächt: ME gilt zwischen \textit{Schreiber–Schreiber} und \textit{Schreiber–Leser}, nicht zwischen \textit{Leser–Leser}.


b) Implementierung mit \texttt{mutex = 1} (binär), \texttt{wrt = 1} (binär), \texttt{readers = 0}:


\begin{lstlisting}
acquire_read():
    P(mutex)
    readers++
    if (readers == 1) P(wrt)   // erster Leser sperrt Schreiber aus
    V(mutex)

release_read():
    P(mutex)
    readers--
    if (readers == 0) V(wrt)   // letzter Leser gibt frei
    V(mutex)

acquire_write():
    P(wrt)

release_write():
    V(wrt)
\end{lstlisting}


Begründung: \texttt{wrt} wird vom \textit{ersten} Leser angefordert und vom \textit{letzten} freigegeben. Solange mindestens ein Leser drin ist, hält die Lesergruppe \texttt{wrt}, und kein Schreiber kommt durch. \texttt{mutex} schützt nur den Zähler \texttt{readers} selbst.


c) \textbf{Verhungerungsszenario} (Reader-Preference):



{\small
\begin{tabularx}{\linewidth}{XXXX}
\hline
\rowcolor{tableheadbg}
{\color{white}\textbf{Zeit}} & {\color{white}\textbf{Aktion}} & {\color{white}\textbf{readers}} & {\color{white}\textbf{wrt-Status}} \\[2pt]
\hline
\endhead
\hline
\endfoot
t1 & R1 acquire\_read & 1 & belegt (von R1) \\
t2 & W1 acquire\_write — wartet auf wrt & 1 & belegt \\
t3 & R2 acquire\_read (mutex frei → readers=2) & 2 & weiterhin belegt \\
t4 & R3 acquire\_read & 3 & belegt \\
t5 & R1 release\_read & 2 & belegt \\
t6 & R4 acquire\_read & 3 & belegt \\
… & beliebig viele Leser kommen und gehen, \texttt{readers} sinkt nie auf 0 & ≥1 & belegt \\
\end{tabularx}
}


Solange ständig neue Leser eintreffen, bevor \texttt{readers} auf 0 sinkt, wird \texttt{V(wrt)} nie ausgeführt — W1 wartet endlos. Verletzt ist die \textbf{Fairness-Bedingung} (Bedingung 4): jeder Wartende muss \textit{irgendwann} eintreten dürfen.


\textbf{Behebung (Skizze)}: Eine zusätzliche Vorrang-Sperre einführen, die neue Leser blockiert, sobald ein Schreiber wartet — z. B. eine \texttt{turnstile}-Semaphore, die jeder Leser kurz \texttt{P()}/\texttt{V()}t, und die ein wartender Schreiber zu Beginn von \texttt{acquire\_write} \texttt{P()}t. Dadurch läuft jeder neue Leser in die Sperre, bis der Schreiber fertig ist (Writer-Preference). Alternativ: FIFO-Warteschlange statt Semaphor — Leser und Schreiber werden in Eintreffreihenfolge bedient.





% ──────────────────────────────────────────────────────────────

\end{document}
